% Nouveau EGA V (cote 25): archive cote 25, site page 44 (PDF page 45).
% typescript. Draft compared: none.
% AI-transcribed by Claude Opus 5.5 (claude-opus-5-5, adaptive thinking, effort high), Message Batches API, batch ega25b (msgbatch_01A7n5YuXubUPNK1zsepi8RG), 2026-11-13; not yet checked by eye.
\scanpage{25/44}
\mspage{\hand{40}}

dans T tel que \hand{$\xi$} $\in U(k)$  implique qu'il existe exactement  un point  y de
$\underline{Y}^{\text{sing}}(k)$ au dessus de \hand{$\xi$}. Cela résultera de 8.7. c) qui implique plus
 précisément:

\uline{Corollaire} 8.9.  Supposons la condition (i) de 8.8. satisfaite, et
soit  
U l'ouvert de T des points en lesquels T est lisse sur k. Alors $U\neq 0$, 
$\underline{Y}^{\text{sing}}|U \longrightarrow U$ est une immersion ouverte,  injection $Y^{\text{sing}}|U$ ne contient pas des points de $\underline{Y}^{\unsure{\text{sing sing}}}$.
 Si X est propre sur
k donc $g:\underline{Y}^{\text{sing}}\longrightarrow T$ surjectif, alors $\underline{Y}^{\text{sing}}|U\longrightarrow U$ est un isomorphi
me, et U est le plus grand ouvert de T ayant cette dernière propriété.

Tout d'abord, comme g est dominant et gén  étale, il est
gén. étale, donc on peut \unsure{trouver|tmouver} un ouvert non vide V de T tel que

$\underline{Y}^{\text{sing}}|V \longrightarrow V$ soit étale et surjectif\placedmargin{ce qui implique V lisse$/k$}. Si alors \hand{$\xi$}$\in V(k)$ et si y est
un point de $\underline{Y}^{\text{sing}}(k)$ au dessus de \hand{$\xi$}, alors avec les notations de
8.7. c)  l'espace H' n'est autre que l'espace tangent à T en \hand{$\xi$},
et comme on a noté à cet endroit, cela implique que le point x de X(k)
projection de y est déterminé comme le point orthogonal à H', donc est
uniquement déterminé, donc comme $\underline{Y}^{\text{sing}}\subset X\text{x}$\hand{$\check{P}$}, y est uniquement déter-%
miné. Cela prouve déjà que g est birationnel (étant dominant, 
génériquement  étale et gén radiciel). 
D'autre part, le morphisme \hand{$\unsure{\nu|\gamma}$}(dont la définition en forme est évidente)
qui associe à tout $\xi$$\in U(k)$ l'unique point $x=$ \hand{$\nu$}(\hand{$\xi$}) $\in P$ orthogonal à
l'espace tangent à U en \hand{$\xi$}, coïncide sur V avec le composé
$V\rightarrow \underline{Y}^{\text{sing}}|V \longrightarrow X$, où la deuxième flèche est la projection; donc
posant  $h=($\hand{$\nu$}$,\mathrm{id}) : U\rightarrow P\text{x}T$ , 
 $g_1 = g|g^{-1}(U)$: $g^{-1}(U)\rightarrow U$ , le composé

\placedmargin{44}
